\subsection{Formal Hypotheses}

Based on the framework above, we derive the following testable hypotheses.

\paragraph{H1 (Polarization and Absolute Price Reactions).}
Political polarization increases the dispersion of interpretations of auditor changes, leading to larger absolute abnormal returns around the disclosure.

\paragraph{H2 (Polarization and Trading Volume).}
Political polarization increases disagreement-driven trading, resulting in higher abnormal trading volume around auditor change announcements.

\paragraph{H3 (Signal Ambiguity).}
The effects of political polarization on price reactions and trading volume are stronger for more ambiguous auditor changes, such as dismissals, switches involving lower-reputation auditors, or disclosures with indications of disagreement.

\paragraph{H4 (Investor Base Mechanism).}
The effects of political polarization are stronger for firms with more ideologically dispersed investor bases, consistent with polarization operating through heterogeneous belief formation across investors.

\subsection{Mapping Theory to Empirical Design}

The theoretical framework maps directly into the empirical specification. Absolute abnormal returns proxy for dispersion in posterior beliefs, while abnormal trading volume captures the intensity of disagreement-driven trading.

Auditor change events provide sharp information shocks. Cross-sectional variation in political polarization captures differences in how investors process these shocks. The interaction between auditor change indicators and polarization measures identifies whether polarization amplifies the market consequences of audit-related disclosures.

Heterogeneity tests across types of auditor changes provide evidence on the role of signal ambiguity. Investor-base polarization measures provide direct tests of the mechanism.

\subsection{Implications for Informational Efficiency}

A key implication of the framework is that political polarization affects the informational efficiency of audit-related disclosures. Even if auditor changes convey economically meaningful information, polarization can lead to persistent disagreement and delayed convergence of beliefs.

This implies that the effectiveness of auditing as a credibility-enhancing institution depends not only on the quality of the audit, but also on the interpretive environment in which audit signals are received.

\subsection{Positioning and Contribution}

This paper contributes to three literatures.

First, it extends the audit and capital markets literature by showing that the pricing of audit signals depends on investor polarization, not just information content.

Second, it contributes to the political economy of finance by demonstrating that polarization affects the interpretation of formal information disclosures.

Third, it contributes to the disagreement literature by identifying a novel source of belief heterogeneity rooted in identity-driven interpretation.

\subsection{Integration with the Full Paper}

To ensure coherence at submission level, the remainder of the paper should align with this framework as follows:

\begin{itemize}
    \item \textbf{Introduction:} Emphasize polarization as an information-processing friction affecting audit signal interpretation.
    \item \textbf{Literature Review:} Explicitly distinguish this paper from CAM studies, auditor change literature, and political trading papers.
    \item \textbf{Empirical Section:} Provide precise definitions of variables, event windows, and polarization measures.
    \item \textbf{Results Section:} Organize results around baseline effects, heterogeneity, mechanism tests, and robustness.
\end{itemize}